\chapter{Einleitung}
\label{ch:einleitung}

\section{Problemstellung und Motivation}
% Ersetzen: Warum ist das Problem relevant? Wo ist die Forschungslücke?
Ein Fadenpendel ist ein einfaches physikalisches System, an dem sich
die Beziehung zwischen Modell und Beobachtung untersuchen lässt.
Die folgende Beispieldarstellung demonstriert die Struktur einer
wissenschaftlichen Arbeit, stellt aber selbst keine Forschungsleistung dar.

\section{Forschungsfrage und Hypothese}
% Die Forschungsfrage muss präzise und messbar bzw. prüfbar sein.
\textbf{Mögliche Forschungsfrage:} Wie hängt die Schwingungsdauer eines
Fadenpendels bei kleinen Auslenkwinkeln von der Pendellänge ab?

\textbf{Hypothese:} Unter idealisierten Bedingungen wächst die
Schwingungsdauer proportional zur Quadratwurzel der Pendellänge.

\section{Ziel und Aufbau der Arbeit}
In Kapitel~\ref{ch:theorie} werden die theoretischen Grundlagen
beschrieben, in Kapitel~\ref{ch:methodik} ein möglicher Versuchsablauf
skizziert. Kapitel~\ref{ch:ergebnisse} zeigt ausschließlich
\emph{berechnete Modellwerte}; ihre Grenzen werden in
Kapitel~\ref{ch:diskussion} diskutiert.
